\documentclass[11pt,oneside]{report}

\usepackage[margin=0.95in,headheight=24pt]{geometry}
\usepackage[T1]{fontenc}
\usepackage[utf8]{inputenc}
\usepackage{amsmath,amssymb}
\usepackage{booktabs}
\usepackage{array}
\usepackage{xcolor}
\usepackage{hyperref}
\usepackage{enumitem}
\usepackage{listings}
\usepackage{titlesec}
\usepackage{fancyhdr}
\usepackage[protrusion=true,expansion=false]{microtype}

\definecolor{Ink}{HTML}{1B2430}
\definecolor{Muted}{HTML}{5B6472}
\definecolor{Accent}{HTML}{0077B6}
\definecolor{AccentTwo}{HTML}{D95F02}
\definecolor{Soft}{HTML}{F5F7FA}

\renewcommand{\familydefault}{\sfdefault}
\setlength{\parindent}{0pt}
\setlength{\parskip}{6pt}
\setlist{itemsep=2pt,topsep=4pt}

\hypersetup{colorlinks=true,linkcolor=Accent,urlcolor=Accent,citecolor=Accent}

\lstdefinestyle{shellstyle}{
  basicstyle=\ttfamily\small,
  backgroundcolor=\color{Soft},
  columns=fullflexible,
  keepspaces=true,
  breaklines=true,
  frame=leftline,
  framerule=2pt,
  rulecolor=\color{Accent},
  xleftmargin=0.8em,
  framesep=0.6em,
  showstringspaces=false
}
\lstset{style=shellstyle}
\newcommand{\code}[1]{\texttt{\small #1}}

\titleformat{\chapter}[block]
  {\huge\bfseries\color{Ink}}{\color{Accent}\thechapter}{0.75em}{}
\titlespacing*{\chapter}{0pt}{-0.15in}{1.1em}
\titleformat{\section}{\Large\bfseries\color{Ink}}{\thesection}{0.6em}{}
\titleformat{\subsection}{\large\bfseries\color{Ink}}{\thesubsection}{0.6em}{}

\pagestyle{fancy}
\fancyhf{}
\lhead{\textcolor{Muted}{QuadrupolarDFT}}
\rhead{\textcolor{Muted}{User Manual}}
\cfoot{\textcolor{Muted}{\thepage}}
\renewcommand{\headrulewidth}{0.3pt}

\newcommand{\Cq}{C_{\mathrm{Q}}}
\newcommand{\Vzz}{V_{zz}}

\begin{document}
\pagenumbering{gobble}
\begin{titlepage}
\pagecolor{Soft}
\color{Ink}
\vspace*{1.1in}
{\Huge\bfseries QuadrupolarDFT\par}
\vspace{0.25in}
{\Large User Manual\par}
\vspace{0.12in}
{\large Ab initio electric-field gradients and finite-temperature NQR parameters\par}
\vspace{0.45in}
{\color{Accent}\rule{0.82\linewidth}{3pt}\par}
\vspace{0.35in}
\begin{minipage}{0.80\linewidth}
\large
QuadrupolarDFT is the electric-field-gradient (EFG) workspace within the
MRSpinDynamics repository. It converts first-principles EFG tensors into nuclear
quadrupolar coupling constants, asymmetry parameters, and zero-field NQR
transition frequencies, and corrects those 0\,K static quantities to a finite
temperature through harmonic vibrational averaging driven by a finite-displacement
ABINIT DFPT workflow.
\end{minipage}
\vfill
{\large MRSpinDynamics project \textbullet{} June 2026\par}
\end{titlepage}
\nopagecolor
\clearpage
\pagenumbering{roman}
\tableofcontents
\clearpage
\pagenumbering{arabic}

% =====================================================================
\chapter{Scope and Physical Background}

QuadrupolarDFT addresses one link in the magnetic-resonance modelling chain: it
turns an \emph{ab initio} electronic-structure result into the parameters that
control a nuclear quadrupole resonance (NQR) experiment.

A nucleus with spin $I \geq 1$ carries an electric quadrupole moment $Q$ that
couples to the electric-field-gradient (EFG) tensor $V_{ij} = \partial^2
V/\partial x_i \partial x_j$ produced by the surrounding charge distribution. The
EFG is symmetric and, away from the nucleus, traceless (Laplace's equation), so
in its principal-axis system it is described by two numbers: the largest-magnitude
component $\Vzz$ and the asymmetry parameter
\begin{equation}
\eta = \frac{V_{xx} - V_{yy}}{\Vzz}, \qquad |\Vzz| \geq |V_{yy}| \geq |V_{xx}|,
\qquad 0 \leq \eta \leq 1 .
\end{equation}
The quadrupolar coupling constant follows from $\Vzz$ and the nuclear quadrupole
moment,
\begin{equation}
\Cq = \frac{e\,Q\,\Vzz}{h},
\label{eq:cq}
\end{equation}
and the zero-field NQR transition frequencies follow from $\Cq$, $\eta$ and the
spin $I$ by diagonalizing the quadrupole Hamiltonian
\begin{equation}
\frac{H_{\mathrm Q}}{h} = \frac{\Cq}{4 I (2I-1)}
  \left[\,3 I_z^2 - I(I+1) + \eta\,(I_x^2 - I_y^2)\,\right].
\label{eq:hq}
\end{equation}
For spin $I=1$ (e.g.\ $^{14}$N) this gives the familiar triplet
\begin{equation}
\nu_{\pm} = \tfrac{3}{4}\Cq\left(1 \pm \tfrac{\eta}{3}\right), \qquad
\nu_{0} = \nu_{+} - \nu_{-} = \tfrac{1}{2}\Cq\,\eta .
\label{eq:spin1}
\end{equation}

\section{Why a finite-temperature correction is needed}

A static DFT calculation evaluates $V_{ij}$ at a single, fixed nuclear
configuration at $0$\,K. Measured NQR frequencies, however, are taken at finite
temperature and are strongly temperature dependent: in molecular crystals the
lines typically \emph{fall} as temperature rises (the Bayer law). Two physically
distinct effects separate the static DFT number from the measurement:

\begin{enumerate}
  \item \textbf{Thermal expansion / geometry.} The lattice constants and internal
    coordinates at temperature $T$ differ from the relaxed $0$\,K cell. This is a
    static effect: it is captured by evaluating the EFG on a
    temperature-appropriate structure.
  \item \textbf{Vibrational averaging.} The nucleus samples a thermal
    distribution of positions through the phonons. The measured EFG is the tensor
    averaged over that motion, including zero-point motion, which is present even
    at $0$\,K. For molecular crystals this dynamic term usually dominates:
    low-frequency librations reorient the EFG principal axes and reduce the line
    frequencies.
\end{enumerate}

QuadrupolarDFT provides the static layer (Chapter~\ref{ch:static}) and a harmonic
treatment of the vibrational term (Chapters~\ref{ch:finiteT}--\ref{ch:workflow}).

% =====================================================================
\chapter{Installation and Conventions}

\section{Installation}

The package depends only on NumPy. From the \code{QuadrupolarDFT/} directory:
\begin{lstlisting}
python -m pip install -e .
python -m pip install -e ".[dev]"     # adds ruff for linting
python -m unittest discover -s tests
\end{lstlisting}
On a system whose Python is externally managed (PEP~668), the package can also be
used without installation by setting \code{PYTHONPATH} to the \code{src/}
directory, since NumPy is the only requirement:
\begin{lstlisting}
export PYTHONPATH="$PWD/src"
python3 examples/abinit/efg_temperature.py --help
\end{lstlisting}

\section{Tensor and unit conventions}

\code{EFGTensor} stores the symmetric tensor in SI units (V/m$^2$) and sorts the
principal components as $|\Vzz| \geq |V_{yy}| \geq |V_{xx}|$. It accepts input in
ABINIT atomic units (the default), in SI, or in the $10^{21}$\,V/m$^2$ scale, and
exposes \code{vzz\_si} and \code{eta}. Tensor averaging
(\code{average\_tensors}) operates on the tensor components and only then
diagonalizes the result, so that asymmetry information survives the average.

\section{Link to the spin-dynamics simulator}

The companion \code{spin\_dynamics} package parameterizes a quadrupolar site by
its $\eta=0$ transition frequency $\nu_{\mathrm Q}$. This is related to $\Cq$ by
the prefactor of Eq.~\eqref{eq:hq}:
\begin{equation}
\nu_{\mathrm Q} = \Cq \cdot \frac{d}{4 I (2I-1)} =
\begin{cases}
\tfrac{3}{4}\,\Cq & I = 1, \\[4pt]
\tfrac{1}{2}\,\Cq & I = \tfrac{3}{2},
\end{cases}
\end{equation}
with $d = 3$ for spin~1 and $d = 6$ for spin~3/2. The cross-project
\code{mr\_integration} layer uses this mapping to feed DFT parameters into the
simulator and to compare predictions against the measured NQR database.

% =====================================================================
\chapter{Static EFG Workflow}
\label{ch:static}

The static path produces $\Cq$, $\eta$ and the NQR lines from an ABINIT EFG
calculation.

\begin{enumerate}
  \item \code{format\_abinit\_efg\_block} emits the \code{nucefg}/\code{quadmom}
    input lines that request an EFG calculation (PAW datasets are required;
    norm-conserving pseudopotentials are not suitable).
  \item ABINIT is run on the input (see \code{examples/abinit/run\_nano2\_efg\_wsl.sh}).
  \item \code{parse\_abinit\_efg} parses the \code{Electric Field Gradient
    Calculation} blocks of the output into \code{AbinitEFGRecord} objects, one per
    atom, carrying $\Cq$, $\eta$, the principal values/axes, and the total EFG
    tensor.
  \item \code{coupling\_constant\_hz} and \code{nqr\_frequencies\_hz} convert a
    tensor and a nuclear quadrupole moment into $\Cq$ and the unique zero-field
    transition frequencies (Eqs.~\eqref{eq:cq}--\eqref{eq:hq}).
\end{enumerate}

\code{examples/parse\_abinit\_efg.py} demonstrates parsing an output file. The EFG
components, not the total energy, must be converged against \code{ecut},
\code{pawecutdg} and the $k$-point mesh.

% =====================================================================
\chapter{Finite-Temperature EFG}
\label{ch:finiteT}

\section{Harmonic vibrational averaging}

Expanding the EFG tensor to second order in the mass-weighted phonon normal
coordinates $Q_k$ and averaging over the harmonic distribution gives
\begin{equation}
\langle V_{ij}\rangle(T) = V_{ij}^{\mathrm{eq}}
  + \tfrac{1}{2}\sum_k \frac{\partial^2 V_{ij}}{\partial Q_k^2}\,
  \langle Q_k^2\rangle(T),
\label{eq:avg}
\end{equation}
where the mean-square amplitude of mode $k$ (angular frequency $\omega_k$) is
\begin{equation}
\langle Q_k^2\rangle(T) = \frac{\hbar}{2\omega_k}
  \coth\!\left(\frac{\hbar\omega_k}{2 k_{\mathrm B} T}\right).
\label{eq:q2}
\end{equation}
The $\coth$ factor equals $2 n(\omega_k,T) + 1$ with $n$ the Bose--Einstein
occupation. It tends to $1$ as $T\to 0$ (pure zero-point motion --- which is why
the static EFG is never exactly what is measured) and to
$2 k_{\mathrm B} T / \hbar\omega_k$ in the classical high-temperature limit. These
quantities live in \code{quadrupolar\_dft.thermal}.

\textbf{Tensor-frame averaging.} The average in Eq.~\eqref{eq:avg} is taken on the
tensor components in a fixed crystal frame; the result is diagonalized only
afterwards. Librational motion reorients the principal-axis system, so averaging
$\Vzz$ or $\eta$ directly would discard exactly that effect. As a consequence
$\eta$ is itself temperature dependent. \code{thermally\_averaged\_efg} performs
the average and removes any residual trace (the EFG is traceless by Laplace's
equation; the finite-difference curvatures of Section~\ref{sec:fd} amplify
rounding noise into a spurious trace, which is projected out).
\code{efg\_temperature\_sweep} returns $\Cq(T)$, $\eta(T)$ and the lines at each
requested temperature.

\section{The analytic Bayer limit}

Where per-mode EFG curvatures are not available, the temperature \emph{form} can
be validated directly against a measured line series with the single-libration
Bayer--Kushida model
\begin{equation}
\nu(T) = \nu_0\left[1 - a\,\coth\!\left(\frac{\hbar\omega}{2 k_{\mathrm B}T}\right)\right],
\label{eq:bayer}
\end{equation}
where $a$ is a small dimensionless librational amplitude.
\code{fit\_bayer\_single\_mode} fits $(\nu_0, a, \omega)$ to measured $\nu(T)$
data with a NumPy-only routine: for a fixed $\omega$ the model is linear in
$\nu_0$ and $\nu_0 a$, so a grid scan over $\omega$ with a linear least-squares
solve at each point suffices.

% =====================================================================
\chapter{The DFPT Finite-Displacement Workflow}
\label{ch:workflow}

The mode curvatures $\partial^2 V_{ij}/\partial Q_k^2$ in Eq.~\eqref{eq:avg} come
from displacing the structure along each phonon mode and recomputing the EFG.
Because the DFT runs are expensive and run outside the package, the workflow
separates into three stages with a calculation between each. The CLI
\code{examples/abinit/efg\_temperature.py} drives all three.

\section{Normal-coordinate convention}
\label{sec:fd}

A phonon eigenvector $\boldsymbol\varepsilon_k$ is mass-weighted and normalized
($\sum |\varepsilon_k|^2 = 1$). A normal-coordinate step $\delta Q$ displaces atom
$i$ (mass $m_i$) by
\begin{equation}
\Delta \mathbf r_i = \delta Q\, \frac{\boldsymbol\varepsilon_{k,i}}{\sqrt{m_i}},
\end{equation}
and the curvature is the central difference of the EFG tensors at $0, \pm\delta Q$,
\begin{equation}
\frac{\partial^2 V_{ij}}{\partial Q_k^2} \approx
  \frac{V_{ij}(+\delta Q) - 2 V_{ij}(0) + V_{ij}(-\delta Q)}{\delta Q^2}.
\end{equation}
These units are consistent with $\langle Q_k^2\rangle$ from Eq.~\eqref{eq:q2}, so
the product in Eq.~\eqref{eq:avg} is an EFG.
\code{normal\_coordinate\_step\_si} chooses $\delta Q$ from a target maximum
Cartesian displacement (default $0.05$\,\AA).

\section{Stage 1 --- phonons}

\code{efg\_temperature.py phonon} writes a two-dataset ABINIT input (ground state
plus a $\Gamma$-point DFPT response) and the matching anaddb input, derived from a
converged static EFG input. \code{run\_phonon\_wsl.sh} runs ABINIT, selects the
DFPT derivative database (the largest \code{*\_DDB}), and runs anaddb to produce
the phonon frequencies and eigenvectors.

\begin{lstlisting}
python3 examples/abinit/efg_temperature.py phonon \
    --base examples/abinit/nano2_efg.abi --out runs/nano2_ph
bash examples/abinit/run_phonon_wsl.sh runs/nano2_ph
\end{lstlisting}

The generated DFPT input is a starting template. Two settings are required for
PAW~+~GGA phonons and are emitted automatically: \code{iscf2 7} (the response
driver requires \code{iscf}~$\leq 9$) and \code{pawxcdev 0} (the GGA exchange
correlation kernel is not implemented in the development scheme). Converge the
tolerances and $k$-mesh for production use.

\section{Stage 2 --- displaced EFG runs}

\code{efg\_temperature.py displace} parses the anaddb eigenvectors
(\code{parse\_anaddb\_modes}; acoustic and imaginary modes are dropped, and modes
are paired with frequencies by mode number so that gaps in the eigenvector output
do not misalign them), generates one displaced EFG input per $\pm$ mode plus a
manifest, and writes positions as fractional \code{xred} coordinates.
\code{run\_finite\_displacement\_wsl.sh} runs ABINIT on every input.

\begin{lstlisting}
python3 examples/abinit/efg_temperature.py displace \
    --base examples/abinit/nano2_efg.abi --anaddb runs/nano2_ph/anaddb.out \
    --target 2 --max-modes 6 --out runs/nano2_disp
bash examples/abinit/run_finite_displacement_wsl.sh runs/nano2_disp
\end{lstlisting}

\code{--target} is the $0$-based index of the resonant nucleus; \code{--max-modes}
restricts the calculation to the lowest-frequency librations, which dominate the
temperature dependence. If the anaddb eigenvector layout differs from the parser's
expectation, \code{--modes modes.json} supplies frequencies and eigenvectors
directly.

\section{Stage 3 --- collect}

\code{efg\_temperature.py collect} parses the displaced EFG outputs, central
differences the mode curvatures, runs the temperature sweep, and reports
$\Cq(T)$, $\eta(T)$, the lines, and $\mathrm d\nu/\mathrm dT$.

\begin{lstlisting}
python3 examples/abinit/efg_temperature.py collect \
    --workdir runs/nano2_disp --temperatures 0,77,150,300 --quadmom 0.02044
\end{lstlisting}

% =====================================================================
\chapter{Worked Example: NaNO$_2$ $^{14}$N}

Sodium nitrite is the reference case for the workspace: its $^{14}$N parameters
appear in the DFT runs, the spin-dynamics simulator, and the measured NQR
database. A six-mode finite-displacement calculation on the starter geometry
produces the temperature dependence in Table~\ref{tab:nano2}.

\begin{table}[h]
\centering
\begin{tabular}{rrrl}
\toprule
$T$ (K) & $\Cq$ (MHz) & $\eta$ & lines (MHz) \\
\midrule
static & $-5.173$ & $0.043$ & 0.11, 3.82, 3.94 \\
0   & $-4.963$ & $0.016$ & 0.04, 3.70, 3.74 \\
77  & $-4.927$ & $0.014$ & 0.03, 3.68, 3.71 \\
150 & $-4.816$ & $0.002$ & 0.00, 3.61, 3.61 \\
300 & $-4.534$ & $0.033$ & 0.08, 3.36, 3.44 \\
\bottomrule
\end{tabular}
\caption{Finite-temperature $^{14}$N EFG of NaNO$_2$ from a six-mode harmonic
average on the starter geometry. The slope of the two strong lines is
$\mathrm d\nu/\mathrm dT \approx -1.0$ to $-1.1$\,kHz/K over $0$--$300$\,K.}
\label{tab:nano2}
\end{table}

The result reproduces the expected physics with no fitting: $|\Cq|$ and the lines
decrease smoothly with temperature (the Bayer sign), the $0$\,K value already lies
below the static one by about $4\%$ (zero-point motion), and $\eta$ shifts with
temperature. The slope, $\sim -1.0$\,kHz/K, has the correct sign and order of
magnitude relative to the measured coefficients of $-1.3$\,kHz/K ($\nu_-$) and
$-2.2$\,kHz/K ($\nu_+$).

\textbf{Caveat.} The starter geometry is unrelaxed: its static $\eta = 0.04$
disagrees with the experimental $0.38$, the two strong lines come out
near-degenerate ($\sim 3.7$\,MHz) rather than well split ($3.60$/$4.64$\,MHz), and
the phonon spectrum contains imaginary modes. The pipeline is therefore validated
as physically sensible, but quantitative agreement requires relaxing the structure
to an energy minimum first.

% =====================================================================
\chapter{Validation Against the Measured Database}

The cross-project \code{mr\_integration} layer closes the loop against the
NQRDatabase. \code{measured\_temperature\_coefficients} reads the stored
$\mathrm d\nu/\mathrm dT$ for a compound and isotope;
\code{compare\_temperature\_coefficients} matches a predicted set of
$(\text{frequency}, \mathrm d\nu/\mathrm dT)$ pairs --- from either a
\code{efg\_temperature\_sweep} (via \code{slopes\_from\_temperature\_points}) or a
Bayer fit --- to the measured coefficients by frequency, so that a DFT temperature
dependence can be checked against experiment. The same layer also validates the
static $(\Cq,\eta)$ against measured line positions and supplies the
$\Cq \leftrightarrow \nu_{\mathrm Q}$ conversion used by the simulator.

% =====================================================================
\chapter{Public API Reference}

\section{quadrupolar\_dft.tensors / quadrupolar / abinit}
\begin{itemize}
  \item \code{EFGTensor}, \code{average\_tensors} --- tensor conventions and
    frame-correct averaging.
  \item \code{coupling\_constant\_hz}, \code{nqr\_frequencies\_hz} --- $\Cq$ and
    zero-field transition frequencies.
  \item \code{AbinitEFGRecord}, \code{parse\_abinit\_efg},
    \code{format\_abinit\_efg\_block} --- ABINIT EFG input/output.
\end{itemize}

\section{quadrupolar\_dft.thermal}
\begin{itemize}
  \item \code{thermal\_quantum\_factor}, \code{bose\_occupation},
    \code{mean\_square\_normal\_coordinate}, \code{wavenumber\_to\_angular\_frequency}.
\end{itemize}

\section{quadrupolar\_dft.vibrational}
\begin{itemize}
  \item \code{VibrationalMode}, \code{efg\_curvature\_central\_difference},
    \code{thermally\_averaged\_efg}, \code{efg\_temperature\_sweep},
    \code{ThermalEFGPoint}.
  \item \code{bayer\_frequency}, \code{bayer\_slope\_hz\_per\_k},
    \code{fit\_bayer\_single\_mode}, \code{BayerFit}.
\end{itemize}

\section{quadrupolar\_dft.finite\_displacement / abinit\_phonon}
\begin{itemize}
  \item \code{Crystal}, \code{PhononMode}, \code{parse\_abinit\_structure},
    \code{generate\_displacement\_jobs}, \code{abinit\_input\_with\_positions},
    \code{write\_jobs}.
  \item \code{collect\_efg\_outputs}, \code{vibrational\_modes\_from\_collected},
    \code{vibrational\_modes\_from\_efg}, \code{efg\_tensor\_from\_record}.
  \item \code{phonon\_dfpt\_input}, \code{anaddb\_input},
    \code{parse\_anaddb\_modes}, \code{modes\_from\_arrays}.
\end{itemize}

% =====================================================================
\chapter{Limitations and Roadmap}

\textbf{Current limitations.}
\begin{itemize}
  \item The $\Cq \leftrightarrow \nu_{\mathrm Q}$ scale and several helpers are
    calibrated for spin~1 and spin~3/2; higher half-integer spins are not yet
    handled by the simulator-facing conversion.
  \item The generated DFPT input is a template; PAW DFPT settings, tolerances and
    $k$-mesh must be checked for production accuracy.
  \item The anaddb eigenvector parser targets ABINIT~9.x text output; a
    \code{--modes} JSON escape hatch covers other layouts.
  \item Finite-difference curvatures use the printed EFG precision; large
    displacement amplitudes or extra output digits improve the second derivative.
\end{itemize}

\textbf{Roadmap.}
\begin{itemize}
  \item A structure-relaxation stage so the geometry sits at an energy minimum
    (removing imaginary modes and correcting $\eta$), the main step toward
    quantitative agreement.
  \item A quasi-harmonic thermal-expansion path and AIMD/PIMD snapshot averaging
    for strongly anharmonic systems such as NaNO$_2$ near its ferroelectric
    transition.
  \item A netCDF / phonopy phonon reader to complement the anaddb text parser.
\end{itemize}

\end{document}
